% Pursuing Stacks (cotes 134-2 to 134-8): archive cote 134-4, site page 50 (PDF page 51), typescript page 138.
% typescript. Draft compared: ps3.tex:2524-2585.
% AI-transcribed by Claude Opus 5.5 (claude-opus-5-5, adaptive thinking, effort high), Message Batches API, batch ps2 (msgbatch_01NVwr1Yuo8qJf5rykeGBXPK), 2026-11-12; not yet checked by eye.
\scanpage{134-4/50}
\mspage{138}

Presumably,  this will yield at the same time an axiomatic caracterization
of \xout{the} $W_{(\mathrm{Cat})}$ namely of the notion of weak equivalence for functors between
% NOTE: the right edge of the scan cuts the end of "between" on this line and of "naive" in the line ending in "since the "naive"
categories.

c) At this point,  we could go on and try and carry through the similar
caracterization for test functors $A\to B^\wedge$ ,  where A and B are both test
categories.  There are also some generalities to develop about``morphisms''
between test categories,  which is ripe too for quite a while and cannot
be pushed off indefinitely  -  here would be the right moment surely.
If the expected ``key result'' for test functors $A\to B^\wedge$ carries through
nicely,  it could presumably be applied \inserted{at once} in order to study general test
functors $A\to M$,  and thus get the clues for cutting out ``the'' natural
notion of a canonic\hand{\unsure{-}}al modelizer,  which ``was in the air'' since the ``naive
question'' of par. 46 (page 95).
\xout{But we could get} \xout{\illegible{1}}

d)However there is another approach to canonical modelizers which is
just appearing,  via the idea (described above) of associating canonically
a notion of\hand{\unsure{``}}weak equivalence\hand{\unsure{''}}$W$ to a \xout{contractibility}\inserted[typed]{homotopy} structure  of type
$\mathrm{Hom}_4$ i.e. to a contractibility structure,  subject possibly to some res-%
trictions.  This ties in,  as explained above,  \over{\unsure{W}}{w}ith a closer look at
the relationship between test categories,  and ``coarse'' contractibility
structures (where all objects are contractible).

It would seem unreasonable to push off a) and b)  any longer now -
so I'll begin with these.  I am hesitant however between c) and d)  -
with a feeling that the la\hand{\unsure{t}}ter approach d) may well turn out to be techni-%
cally the most expedient one.  Both have to be carried tr\over{\unsure{x}}{\unsure{h}}ough any\over{\unsure{h}}{\unsure{y}}how ,
and the two together should give a rather accurate picture of what canoni-%
cal modelizers are about.

On the other hand,  there is still quite a bunch of questions which
have been waiting for investigati\hand{on} - \over{\unsure{f}}{f}or instance the list of questions
of nearly three weeks ago (page 42)\over{,}{.}\xout{\unsure{x}} \over{\unsure{a}}{A}mong the six questions stated
there,  four have been settled\hand{,} or are about to be settled through the
previous program (if it works out),  questions 4) and 6) remain,  the
first one \inserted{being} about $A\hand{^\wedge}$\hand{/}being a closed model category,  and \inserted{about} the homot\hand{o}py struc-%
ture of (Cat).  There \over{\unsure{is}}{ar}e a number of more technical questions too,  for
instance I did not finish yet my review of the ``standart'' test categories,
and never wrote down the proof that $\hand{\Delta}^{\mathrm{f}}$ (simplices with face operations
and no degeneracies)  is indeed a weak test category.  But for the time
being all these questions appear as \inserted{somewhat} marginal with respect to the strong
focus the reflection has been gradually taking nearly since the very

\margin{50}
